\section{Coin Change (Minimum Coins)}
\label{sec:dp-coin-change}

\noindent We now begin the chapter's \textbf{unbounded knapsack} family:
every problem from here on lets each item (coin denomination, rod
length) be used an \textbf{unlimited} number of times, which means the
pick branch stays at the same index instead of advancing -- exactly the
recursion shape the Recursion chapter's Combination Sum section
already derived.

\problemstatement{Given an array $\textit{coins}$ of denominations
(each available in \textbf{unlimited} supply) and a target
$\textit{amount}$, find the \textbf{minimum} number of coins needed to
make exactly $\textit{amount}$, or $-1$ if it is impossible.}

\begin{patternbox}
\emph{``Unlimited supply of each denomination, minimize the count
used''} is the unbounded-knapsack analogue of
Section~\ref{sec:dp-subset-sum-target}'s boolean reachability table: the
pick branch stays at the same index (since the same coin may be reused),
and the combination operator is a \textbf{minimum} of counts instead of a
logical OR of reachability.
\end{patternbox}

\subsection*{Deriving the recurrence}

Let $f(\textit{index}, \textit{amount})$ be the minimum number of coins
(using denominations $0..\textit{index}$) needed to make
$\textit{amount}$. Either we exclude denomination $\textit{index}$
entirely ($f(\textit{index}-1, \textit{amount})$), or we use it at least
once (paying $1$ coin, then recursing on the \emph{same} index since it
remains available, with the reduced amount):
\begin{gather*}
  f(\textit{index}, \textit{amount}) = \min\Big(
  f(\textit{index}{-}1, \textit{amount}),\ \textit{take} \Big), \\
  \textit{take} = \textit{coins}[\textit{index}] \le \textit{amount}\ ?\
  1 + f(\textit{index}, \textit{amount}-\textit{coins}[\textit{index}]) :
  \infty.
\end{gather*}
Base case: $f(0, \textit{amount}) = \textit{amount} /
\textit{coins}[0]$ if $\textit{amount}$ is exactly divisible by
$\textit{coins}[0]$, else $\infty$ (using only the first denomination,
either it divides evenly or it is impossible).

\subsection*{Approach 1: Recursion}

\begin{lstlisting}
#include <bits/stdc++.h>
using namespace std;
const int INF = 1e9;

int coinChangeRecursive(int index, int amount, vector<int>& coins) {
    if (index == 0) {
        return (amount % coins[0] == 0) ? amount / coins[0] : INF;
    }

    int notTake = coinChangeRecursive(index - 1, amount, coins);
    int take = INF;
    if (coins[index] <= amount) {
        take = 1 + coinChangeRecursive(index, amount - coins[index], coins);
    }
    return min(notTake, take);
}

int main() {
    vector<int> coins = {1, 2, 5};
    int n = coins.size(), amount = 11;
    int result = coinChangeRecursive(n - 1, amount, coins);
    cout << (result >= INF ? -1 : result) << endl; // 3
    return 0;
}
\end{lstlisting}

\begin{complexitybox}[Approach 1 Complexity]
\textbf{Time.} Exponential -- each call can recurse on the same index
many times (reducing the amount), giving a recursion tree that grows
combinatorially with $\textit{amount}$.
\textbf{Space.} $O(\textit{amount})$ recursion stack in the worst case
(repeatedly taking the smallest coin).
\end{complexitybox}

\subsection*{Approach 2: Memoization}

\begin{lstlisting}
#include <bits/stdc++.h>
using namespace std;
const int INF = 1e9;

int coinChangeMemo(int index, int amount, vector<int>& coins,
                    vector<vector<int>>& memo) {
    if (index == 0) {
        return (amount % coins[0] == 0) ? amount / coins[0] : INF;
    }
    if (memo[index][amount] != -1) return memo[index][amount];

    int notTake = coinChangeMemo(index - 1, amount, coins, memo);
    int take = INF;
    if (coins[index] <= amount) {
        take = 1 + coinChangeMemo(index, amount - coins[index], coins, memo);
    }
    return memo[index][amount] = min(notTake, take);
}

int main() {
    vector<int> coins = {1, 2, 5};
    int n = coins.size(), amount = 11;
    vector<vector<int>> memo(n, vector<int>(amount + 1, -1));
    int result = coinChangeMemo(n - 1, amount, coins, memo);
    cout << (result >= INF ? -1 : result) << endl; // 3
    return 0;
}
\end{lstlisting}

\begin{complexitybox}[Approach 2 Complexity]
\textbf{Time.} $O(n \cdot \textit{amount})$.
\textbf{Space.} $O(n \cdot \textit{amount})$ memo $+$
$O(\textit{amount})$ recursion stack.
\end{complexitybox}

\subsection*{Approach 3: Tabulation}

\begin{lstlisting}
#include <bits/stdc++.h>
using namespace std;
const int INF = 1e9;

int coinChangeTabulation(vector<int>& coins, int amount) {
    int n = coins.size();
    vector<vector<int>> dp(n, vector<int>(amount + 1, INF));

    for (int a = 0; a <= amount; a++) {
        if (a % coins[0] == 0) dp[0][a] = a / coins[0];
    }

    for (int index = 1; index < n; index++) {
        for (int a = 0; a <= amount; a++) {
            int notTake = dp[index - 1][a];
            int take = INF;
            if (coins[index] <= a) {
                take = 1 + dp[index][a - coins[index]];
            }
            dp[index][a] = min(notTake, take);
        }
    }
    int result = dp[n - 1][amount];
    return result >= INF ? -1 : result;
}

int main() {
    vector<int> coins = {1, 2, 5};
    cout << coinChangeTabulation(coins, 11) << endl; // 3
    return 0;
}
\end{lstlisting}

\subsection*{Dry run of the tabulation table}

\dryrun{Take $\textit{coins} = [1,2,5]$, $\textit{amount}=11$.}

Row $0$ (only coin $1$): $\textit{dp}[0][a] = a$ for all $a$ (every
amount needs exactly $a$ coins of value $1$). Row $1$ (coins $1,2$):
e.g.\ $\textit{dp}[1][11] = \min(\textit{dp}[0][11],\ 1+\textit{dp}[1][9])$;
unrolling, using coin $2$ as many times as possible gives
$\textit{dp}[1][11]=6$ ($5$ twos and one... actually $11=2\times5+1$,
so $5$ coins of $2$ plus $1$ coin of $1$ $=6$ coins). Row $2$ (coins
$1,2,5$): $\textit{dp}[2][11] = \min(\textit{dp}[1][11],\
1+\textit{dp}[2][6])$; following this through, the optimal is
$5+5+1=11$ using $3$ coins, giving $\textit{dp}[2][11]=3$.

\begin{complexitybox}[Approach 3 Complexity]
\textbf{Time.} $O(n \cdot \textit{amount})$.
\textbf{Space.} $O(n \cdot \textit{amount})$ table, no recursion stack.
\end{complexitybox}

\subsection*{Approach 4: Space Optimization}

\begin{lstlisting}
#include <bits/stdc++.h>
using namespace std;
const int INF = 1e9;

int coinChangeSpaceOptimized(vector<int>& coins, int amount) {
    int n = coins.size();
    vector<int> prev(amount + 1, INF);

    for (int a = 0; a <= amount; a++) {
        if (a % coins[0] == 0) prev[a] = a / coins[0];
    }

    for (int index = 1; index < n; index++) {
        vector<int> current(amount + 1, INF);
        for (int a = 0; a <= amount; a++) {
            int notTake = prev[a];
            int take = INF;
            if (coins[index] <= a) {
                take = 1 + current[a - coins[index]]; // same row! (unbounded)
            }
            current[a] = min(notTake, take);
        }
        prev = current;
    }
    int result = prev[amount];
    return result >= INF ? -1 : result;
}

int main() {
    vector<int> coins = {1, 2, 5};
    cout << coinChangeSpaceOptimized(coins, 11) << endl; // 3
    return 0;
}
\end{lstlisting}

\begin{ideabox}[Why ``Take'' Reads from \textit{current}, Not \textit{prev}]
Because the pick branch in an \textbf{unbounded} knapsack recurrence
stays at the \emph{same} index, the space-optimized version must read
the take-branch value from the row currently being built
($\textit{current}$), not from the previous row ($\textit{prev}$) --
the opposite of every 0/1 knapsack-family problem earlier in this
chapter, where both branches read only from $\textit{prev}$. This
single detail is the only change space optimization needs when moving
from the 0/1 family to the unbounded family.
\end{ideabox}

\begin{complexitybox}[Approach 4 Complexity]
\textbf{Time.} $O(n \cdot \textit{amount})$.
\textbf{Space.} $O(\textit{amount})$ for the rolling row.
\end{complexitybox}

\begin{takeawaybox}
Coin Change is 0/1 Knapsack's unbounded twin: same boolean-to-minimum
combination idea, but with the pick branch staying at the same index --
exactly the same 0/1-versus-unbounded distinction the Recursion
chapter drew between Subset Sum II (no reuse) and Combination Sum
(unlimited reuse). The single space-optimization wrinkle above (reading
``take'' from the current row) is the one new implementation detail this
unbounded family introduces.
\end{takeawaybox}
